\documentclass[12pt,a4paper]{article}
\usepackage[margin=25mm]{geometry}
\usepackage[T1]{fontenc}
\usepackage[utf8]{inputenc}
\usepackage{lmodern}
\usepackage{graphicx}
\usepackage{amsmath,amssymb}
\usepackage{booktabs}
\usepackage{tabularx}
\usepackage{float}
\usepackage{subcaption}
\usepackage{circuitikz}
\usepackage{hyperref}
\usepackage{fancyhdr}
\usepackage{titlesec}
\usepackage{listings}
\usepackage{xcolor}

\graphicspath{{Figures/}{images/}}
\hypersetup{colorlinks=true,linkcolor=black,urlcolor=blue}
\pagestyle{fancy}
\fancyhf{}
\lhead{EN3160 -- Image Processing \\ and Machine Vision}
\fancyhead[R]{Intensity Transformations \& \\ Neighborhood Filtering}
\cfoot{\thepage}
\setlength{\headheight}{15pt}
\titleformat{\section}{\Large\bfseries}{\thesection}{0.7em}{}
\titleformat{\subsection}{\large\bfseries}{\thesubsection}{0.7em}{}

% Code listing styling
\definecolor{codegreen}{rgb}{0,0.6,0}
\definecolor{codegray}{rgb}{0.5,0.5,0.5}
\definecolor{codepurple}{rgb}{0.58,0,0.82}
\definecolor{backcolour}{rgb}{0.96,0.96,0.96}

\lstdefinestyle{pythonstyle}{
    backgroundcolor=\color{backcolour},   
    commentstyle=\color{codegreen},
    keywordstyle=\color{magenta},
    numberstyle=\tiny\color{codegray},
    stringstyle=\color{codepurple},
    basicstyle=\ttfamily\footnotesize,
    breakatwhitespace=false,         
    breaklines=true,                 
    captionpos=b,                    
    keepspaces=true,                 
    numbers=left,                    
    numbersep=5pt,                  
    showspaces=false,                
    showstringspaces=false,
    showtabs=false,                  
    tabsize=2
}
\lstset{style=pythonstyle}

\begin{document}

\begin{titlepage}
  \centering
  \vspace*{7mm}
  \includegraphics[width=4.2cm]{campus_logo.png}\par\vspace{7mm}
  {\Large\bfseries Department of Electronic and Telecommunication Engineering\par}
  \vspace{2mm}{\large University of Moratuwa\par}
  \vspace{18mm}{\large B.Sc. Engineering -- Semester 5\par}
  \vspace{5mm}{\Large\bfseries EN3160 -- Image Processing and Machine Vision\par}
  \vspace{15mm}{\Huge\bfseries Intensity Transformations \& Neighborhood Filtering\par} 
  \vfill
  \renewcommand{\arraystretch}{1.35}
  \begin{tabular}{rl}
    \textbf{Laboratory:} & Analog Lab \\
    \textbf{Group number:} & 6 \\
    \textbf{Members:} 
                    & 230449N -- Nuwanaka W. A. S. \\
  \end{tabular}
  \vfill
  {\large \today\par}
\end{titlepage}

\tableofcontents
\newpage

% ---------------------------------------------------------
% QUESTION 1
% ---------------------------------------------------------
\section{Question 1: Piecewise Linear Intensity Transformation}

\subsection{Task Description}
Implement a piecewise linear intensity transformation function \texttt{intensity\_transform(im, breakpoints)} where \texttt{breakpoints} defines a mapping from input intensities $x$ to output intensities $y$, e.g.,
\begin{equation}
\text{breakpoints} = \begin{bmatrix}
0 & 0 \\
50 & 50 \\
100 & 150 \\
150 & 255 \\
255 & 255
\end{bmatrix}
\end{equation}

\subsection{Requirements}
\begin{enumerate}
    \item Implement \texttt{intensity\_transform(im, breakpoints)} using lookup tables (LUT) or piecewise linear interpolation (\texttt{np.interp}).
    \item Load target image \texttt{im01.png} (or \texttt{im01small.png}).
    \item Experiment with different breakpoints to achieve contrast enhancement.
    \item Display the transformation curve plot alongside the original and transformed images side-by-side.
\end{enumerate}

\subsection{Implementation Code Stub}
\begin{lstlisting}[language=Python, caption={Piecewise Linear Intensity Transformation Stub}]
import cv2 as cv
import numpy as np

def intensity_transform(im: np.ndarray, breakpoints: np.ndarray) -> np.ndarray:
    """
    Applies piecewise linear intensity transformation to an input image.
    
    Parameters:
        im (np.ndarray): Input image (grayscale or color)
        breakpoints (np.ndarray): Nx2 matrix of [input_intensity, output_intensity]
        
    Returns:
        np.ndarray: Transformed image of same shape and uint8 dtype
    """
    # TODO: Extract x and y coordinates from breakpoints
    # TODO: Create Look-Up Table (LUT) for intensity mapping [0..255]
    # TODO: Apply cv.LUT or np.interp to transform image intensities
    pass
\end{lstlisting}

\subsection{Discussion \& Observations}
\textit{Discuss how changing breakpoints alters image contrast, highlights, and shadow retention...}

\newpage

% ---------------------------------------------------------
% QUESTION 2
% ---------------------------------------------------------
\section{Question 2: Tissue Accentuation in Brain Proton Density Image}

\subsection{Task Description}
Apply piecewise linear intensity transformations on the brain proton density slice (\texttt{im02.png} / \texttt{im02small.png}) to selectively accentuate:
\begin{itemize}
    \item \textbf{(a)} White matter
    \item \textbf{(b)} Gray matter
\end{itemize}

\subsection{Requirements}
\begin{enumerate}
    \item Analyze histogram of the brain proton density slice to identify intensity ranges corresponding to white matter and gray matter.
    \item Design custom breakpoint curves for (a) White Matter accentuation and (b) Gray Matter accentuation.
    \item Plot transformation curves and display the original vs accentuated images.
\end{enumerate}

\subsection{Implementation Code Stub}
\begin{lstlisting}[language=Python, caption={Tissue Accentuation Code Stub}]
# TODO: Load brain proton density slice 'im02.png' or 'im02small.png'
# TODO: Define breakpoint arrays for white matter and gray matter enhancement
# TODO: Call intensity_transform() for both cases
# TODO: Plot transformation curves and display original, white matter, and gray matter images
\end{lstlisting}

\subsection{Discussion \& Observations}
\textit{Discuss intensity range selections for white and gray matter and visual segmentations achieved...}

\newpage

% ---------------------------------------------------------
% QUESTION 3
% ---------------------------------------------------------
\section{Question 3: Gamma Correction in \texorpdfstring{$L^*a^*b^*$}{L*a*b*} Color Space}

\subsection{Task Description}
Consider the image shown in Fig. 3 (\texttt{im03.png} / \texttt{im03small.png}).
\begin{enumerate}
    \item \textbf{(a)} Convert image to $L^*a^*b^*$ color space and apply gamma correction to the $L^*$ plane. State chosen $\gamma$ value.
    \item \textbf{(b)} Display original vs gamma-corrected image, along with histograms of the $L^*$ plane before and after correction.
\end{enumerate}

\subsection{Mathematical Formulation}
\begin{equation}
L_{\text{out}} = 255 \times \left(\frac{L_{\text{in}}}{255}\right)^\gamma
\end{equation}

\subsection{Implementation Code Stub}
\begin{lstlisting}[language=Python, caption={Gamma Correction Code Stub}]
def gamma_correct_L_plane(im_bgr: np.ndarray, gamma: float) -> tuple[np.ndarray, np.ndarray, np.ndarray]:
    """
    Applies gamma correction to the L* plane in L*a*b* color space.
    """
    # TODO: Convert BGR to Lab space using cv.cvtColor(im_bgr, cv.COLOR_BGR2LAB)
    # TODO: Extract L, a, b planes
    # TODO: Compute normalized power law transform: L_out = 255 * ((L_in / 255) ** gamma)
    # TODO: Merge L_out, a, b and convert back to BGR space
    pass
\end{lstlisting}

\subsection{Discussion \& Observations}
\textit{Explain the choice of gamma value, impact on shadows/highlights in Lab space vs RGB space...}

\newpage

% ---------------------------------------------------------
% QUESTION 4
% ---------------------------------------------------------
\section{Question 4: Vibrance Enhancement in HSV Color Space}

\subsection{Task Description}
Enhance photograph vibrance on image \texttt{im04.jpg} (\texttt{im04small.jpg}) by applying the intensity transformation function:
\begin{equation}
f(x) = \min\left(x + a \times 128 e^{-\frac{(x-128)^2}{2\sigma^2}}, 255\right)
\end{equation}
to the Saturation ($S$) plane in HSV color space, where $x \in [0, 255]$, $a \in [0, 1]$, and $\sigma = 70$.

\subsection{Requirements}
\begin{enumerate}
    \item \textbf{(a)} Split \texttt{im04.jpg} into Hue ($H$), Saturation ($S$), and Value ($V$) planes.
    \item \textbf{(b)} Apply the vibrance transform formula to the $S$ plane.
    \item \textbf{(c)} Adjust $a$ for optimal visual results and report chosen $a$.
    \item \textbf{(d)} Recombine $H, S_{\text{enhanced}}, V$ back into BGR image.
    \item \textbf{(e)} Display original image, vibrance-enhanced image, and transformation plot $f(x)$ vs $x$.
\end{enumerate}

\subsection{Implementation Code Stub}
\begin{lstlisting}[language=Python, caption={Vibrance Enhancement Code Stub}]
def vibrance_transform(S_plane: np.ndarray, a: float, sigma: float = 70.0) -> np.ndarray:
    """Applies vibrance enhancement function to Saturation channel."""
    # TODO: Compute transformation f(x) = min(x + a * 128 * exp(-((x-128)^2)/(2 * sigma^2)), 255)
    # TODO: Cast back to uint8
    pass
\end{lstlisting}

\subsection{Discussion \& Observations}
\textit{Report the chosen 'a' parameter and explain why Gaussian-weighted saturation boosting avoids clipping vibrant colors...}

\newpage

% ---------------------------------------------------------
% QUESTION 5
% ---------------------------------------------------------
\section{Question 5: Custom Histogram Equalization}

\subsection{Task Description}
Write a custom function from scratch to carry out histogram equalization on \texttt{im05.jpg}.

\subsection{Requirements}
\begin{enumerate}
    \item Implement \texttt{custom\_histogram\_equalization(img\_gray)} without using OpenCV's \texttt{cv.equalizeHist}.
    \item Compute intensity frequency histogram $h(r_k)$.
    \item Calculate Probability Density Function (PDF) $p(r_k) = h(r_k) / MN$.
    \item Compute Cumulative Distribution Function (CDF) $T(r_k) = \sum p(r_k)$.
    \item Map intensity levels $s_k = \text{round}(255 \cdot T(r_k))$.
    \item Display image before and after equalization alongside their histograms.
\end{enumerate}

\subsection{Implementation Code Stub}
\begin{lstlisting}[language=Python, caption={Custom Histogram Equalization Stub}]
def custom_histogram_equalization(img_gray: np.ndarray) -> tuple[np.ndarray, np.ndarray, np.ndarray]:
    """Computes histogram equalization manually."""
    # TODO: Calculate intensity counts for values 0 to 255
    # TODO: Compute normalized CDF
    # TODO: Scale CDF to [0..255] and round
    # TODO: Map input image pixels to new intensities
    pass
\end{lstlisting}

\subsection{Discussion \& Observations}
\textit{Compare your manual histogram equalization output with OpenCV cv.equalizeHist and analyze contrast distribution...}

\newpage

% ---------------------------------------------------------
% QUESTION 6
% ---------------------------------------------------------
\section{Question 6: Selective Foreground Histogram Equalization}

\subsection{Task Description}
Apply histogram equalization \textbf{only to the foreground} of an image (e.g. \texttt{q1.jpeg}).

\subsection{Steps}
\begin{enumerate}
    \item \textbf{(a)} Split image into $H, S, V$ planes and display each plane in grayscale.
    \item \textbf{(b)} Select appropriate plane to threshold and extract binary foreground mask $M$.
    \item \textbf{(c)} Extract foreground using \texttt{cv.bitwise\_and} and compute foreground histogram.
    \item \textbf{(d)} Compute cumulative sum of histogram using \texttt{np.cumsum}.
    \item \textbf{(e)} Equalize foreground intensities based on cumulative histogram.
    \item \textbf{(f)} Recombine background with equalized foreground.
\end{enumerate}

\subsection{Implementation Code Stub}
\begin{lstlisting}[language=Python, caption={Selective Foreground Equalization Stub}]
# TODO: Load Fig. 6 image ('q1.jpeg')
# TODO: Convert to HSV and display H, S, V planes
# TODO: Threshold selected plane to create binary foreground mask
# TODO: Apply cv.bitwise_and to extract foreground pixels
# TODO: Compute cumulative sum of foreground histogram using np.cumsum()
# TODO: Equalize foreground intensities and recombine with background
\end{lstlisting}

\subsection{Discussion \& Observations}
\textit{Discuss plane selection choice for segmentation and compare selective foreground equalization vs global equalization...}

\newpage

% ---------------------------------------------------------
% QUESTION 7
% ---------------------------------------------------------
\section{Question 7: Sobel Filtering \& Separable Filtering}

\subsection{Task Description}
Perform Sobel filtering on \texttt{q2.jpeg} (Einstein image) using three distinct methods:
\begin{enumerate}
    \item \textbf{(a)} Using OpenCV \texttt{cv.filter2D} with full 2D Sobel kernel.
    \item \textbf{(b)} Manual implementation of 2D convolution for Sobel filtering.
    \item \textbf{(c)} Utilizing kernel separability property:
    \begin{equation}
    \begin{bmatrix} 1 & 0 & -1 \\ 2 & 0 & -2 \\ 1 & 0 & -1 \end{bmatrix} = \begin{bmatrix} 1 \\ 2 \\ 1 \end{bmatrix} * \begin{bmatrix} 1 & 0 & -1 \end{bmatrix}
    \end{equation}
    Perform sequential 1D horizontal and vertical convolutions.
\end{enumerate}

\subsection{Implementation Code Stub}
\begin{lstlisting}[language=Python, caption={Sobel Filtering Implementations Stub}]
def sobel_filter2d_cv(im: np.ndarray) -> np.ndarray:
    """(a) Sobel filtering using cv.filter2D"""
    pass

def sobel_filter2d_manual(im: np.ndarray) -> np.ndarray:
    """(b) Sobel filtering using manual 2D convolution loop"""
    pass

def sobel_filter_separable(im: np.ndarray) -> np.ndarray:
    """(c) Sobel filtering using 1D separable kernel decomposition"""
    pass
\end{lstlisting}

\subsection{Discussion \& Observations}
\textit{Compare execution times and accuracy of 2D vs 1D separable filtering...}

\newpage

% ---------------------------------------------------------
% QUESTION 8
% ---------------------------------------------------------
\section{Question 8: Image Zooming (Nearest-Neighbor \& Bilinear)}

\subsection{Task Description}
Write a custom program/function \texttt{zoom\_image(im, scale\_factor, method)} to resize images by factor $s \in (0, 10]$.

\subsection{Requirements}
\begin{enumerate}
    \item Implement interpolation methods: (a) Nearest-neighbor interpolation, (b) Bilinear interpolation.
    \item Scale up small images by a factor of 4.
    \item Compute \textbf{Normalized Sum of Squared Differences (SSD)} by comparing upscaled images against original high-resolution targets:
    \begin{equation}
    \text{SSD}_{\text{norm}} = \frac{\sum_{i,j} (I_{\text{orig}}(i,j) - I_{\text{zoomed}}(i,j))^2}{\sum_{i,j} I_{\text{orig}}(i,j)^2}
    \end{equation}
\end{enumerate}

\subsection{Implementation Code Stub}
\begin{lstlisting}[language=Python, caption={Image Zooming Code Stub}]
def zoom_image(im: np.ndarray, factor: float, method: str = 'bilinear') -> np.ndarray:
    """Zooms image by given scale factor using custom interpolation."""
    pass

def compute_normalized_ssd(im_orig: np.ndarray, im_zoomed: np.ndarray) -> float:
    """Computes normalized sum of squared differences."""
    pass
\end{lstlisting}

\subsection{Discussion \& Observations}
\textit{Compare SSD scores for Nearest Neighbor vs Bilinear interpolation and analyze visual artifacts...}

\newpage

% ---------------------------------------------------------
% QUESTION 9
% ---------------------------------------------------------
\section{Question 9: GrabCut Segmentation \& Background Blurring}

\subsection{Task Description}
Consider the Marguerite Daisy flower image.
\begin{enumerate}
    \item \textbf{(a)} Use \texttt{cv.grabCut} to segment flower foreground. Show final segmentation mask, extracted foreground image, and extracted background image.
    \item \textbf{(b)} Produce enhanced image with a substantially blurred background (using Gaussian filter). Display original alongside enhanced image.
    \item \textbf{(c)} \textbf{Discussion Question:} Why is the background area just beyond the edge of the flower quite dark in the enhanced image?
\end{enumerate}

\subsection{Implementation Code Stub}
\begin{lstlisting}[language=Python, caption={GrabCut Segmentation Code Stub}]
# TODO: Load flower image
# TODO: Define bounding box surrounding flower foreground
# TODO: Run cv.grabCut to compute segmentation mask
# TODO: Apply Gaussian Blur to background image and composite
\end{lstlisting}

\subsection{Discussion \& Answer to Part (c)}
\textit{Explain why color bleeding / halo artifacts occur along boundary pixels when applying spatial blurring to background...}

\newpage

% ---------------------------------------------------------
% QUESTION 10
% ---------------------------------------------------------
\section{Question 10: Edge-Preserving Bilateral Filtering}

\subsection{Task Description}
\begin{enumerate}
    \item \textbf{(a)} Use OpenCV \texttt{cv.bilateralFilter} with spatial parameter $\sigma_s$ and range parameter $\sigma_r$. Display original image, Gaussian blurred image, and bilateral filtered image side-by-side.
    \item \textbf{(b)} Implement custom bilateral filter from scratch. Compare results with OpenCV's \texttt{cv.bilateralFilter}. Evaluate results with a quantitative measure (e.g. MSE, PSNR, or SSIM).
\end{enumerate}

\subsection{Implementation Code Stub}
\begin{lstlisting}[language=Python, caption={Custom Bilateral Filter Stub}]
def custom_bilateral_filter(im: np.ndarray, d: int, sigma_color: float, sigma_space: float) -> np.ndarray:
    """Custom implementation of 2D Bilateral Filter."""
    pass
\end{lstlisting}

\subsection{Discussion \& Observations}
\textit{Discuss edge preservation behavior, effect of sigma\_color vs sigma\_space, and quantitative evaluation metrics...}

\end{document}